\begin{abstract}
Unsupervised graph domain adaptation (UGDA) trains without target labels, yet
deployment must also choose an algorithm, configuration, seed, and checkpoint
without those labels. We introduce \benchmark, an auditable trajectory
benchmark that separates this model-selection problem from adaptation
training and seals target labels behind a one-time evaluator. We further study
\method, a covariance-aware risk-recovery module based on an exact identity:
pairwise prediction disagreement equals the sum of individual risks minus
cross-model error covariance. A tight target-graph spectral frame decomposes
this identity without changing total hard-label risk, and our analysis
quantifies how covariance misspecification propagates to recovered risk.
On 44 source-simulated development folds with 675 candidates per task,
covariance correction reduces mean normalized regret from $0.05057$ to
$0.01415$; the additional gain from spectral conditioning over an
equal-information global correction is modest and statistically
inconclusive. On four pre-registered real-target open-development tasks, the
strongest diagnostic pipeline---Agreement top-$20\%$ followed by Transfer
Score---reduces mean normalized regret from $0.216807$ to $0.087507$, but is
not a \method result and fails the registered robustness criteria: worst-task
regret is $0.223881$, only $2/4$ tasks are non-inferior to Transfer Score, and
the source-simulated family-out CVaR guard is unmet. A shortlist-specific
decomposition attributes $79.5\%$ of its residual regret to reranking after
the oracle trajectory is already represented. We therefore stopped
development without freezing a selector or querying the final twelve sealed
transfers. The resulting benchmark, recovery theory, and negative audit
identify shortlist-internal top-ranking calibration---not candidate
coverage---as the central unresolved problem.
\end{abstract}
